\begin{onehalfspace}

Il post-processing del Capitolo~\ref{chap:risultati} recupera candidati ancora presenti nell'istogramma, ma non ripristina informazione distrutta durante l'evoluzione. La \ac{QEC} interviene invece durante il calcolo. Questo capitolo ne introduce gli elementi operativi --- stabilizzatori, sindrome e decoder --- e li verifica sul codice a ripetizione e sul codice di Steane $[[7,1,3]]$; il Capitolo~\ref{chap:surface} passa poi al surface code scalabile.

% -----------------------------------------------
\section{Dalla Ridondanza Ingenua alla Codifica}
\label{sec:ridondanza_codifica}

L'idea di proteggere l'informazione replicandola è classica: si trasmettono tre copie di un bit e si decide a maggioranza. Sul piano quantistico questa strada è sbarrata due volte. Il teorema di no-cloning~\cite{wootters1982} (Sezione~\ref{chap:fondamenti}) impedisce di copiare uno stato ignoto; e la misura diretta di un qubit, necessaria per un ipotetico ``voto di maggioranza'', distrugge la sovrapposizione che si voleva proteggere. Per questo la pratica NISQ più diffusa si limita a una ridondanza \emph{a livello di esecuzione} --- eseguire istanze gemelle dello stesso circuito su regioni diverse del dispositivo e usare l'esito della copia sopravvissuta --- che paga un raddoppio di qubit fisici senza agire sulla causa degli errori.

La \ac{QEC} risolve il problema in modo strutturalmente diverso~\cite{shorcorrection1995, steane1996, roffe2019}: l'informazione di un \textbf{qubit logico} (\textit{logical qubit}) viene \emph{codificata} --- non copiata --- nelle correlazioni di entanglement di più \textbf{qubit fisici} (\textit{physical qubits}), e gli errori vengono diagnosticati misurando operatori ausiliari che rivelano \emph{che cosa è andato storto} senza rivelare --- e quindi senza distruggere --- \emph{che cosa è codificato}. È questa distinzione, non la semplice ridondanza, a rendere possibile la correzione durante il calcolo.

% -----------------------------------------------
\section{Il Ciclo della Correzione: Stabilizzatori, Sindrome, Decoder}
\label{sec:ciclo_qec}

Il funzionamento di ogni codice correttore quantistico segue lo stesso ciclo concettuale:

\begin{enumerate}
    \item \textbf{Codifica}: lo stato logico $\ket{\psi_L}$ viene preparato su $n$ qubit fisici;
    \item \textbf{Rumore}: i canali fisici del Capitolo~\ref{chap:rumore} agiscono sui qubit fisici;
    \item \textbf{Estrazione della sindrome}: qubit ausiliari (\textit{ancilla}) misurano gli stabilizzatori del codice; l'esito --- una stringa di bit classica detta \textbf{sindrome} (\textit{syndrome}) --- identifica la classe di errore senza toccare l'informazione logica;
    \item \textbf{Decodifica}: un algoritmo classico, il \textbf{decoder}, interpreta la sindrome e propone la correzione più probabile;
    \item \textbf{Correzione}: la correzione viene applicata fisicamente, oppure registrata in un registro classico --- il \textbf{Pauli frame} --- che tiene memoria delle correzioni da interpretare sulle misure successive, riducendo le operazioni fisiche e la propagazione di errori.
\end{enumerate}

\paragraph{Stabilizzatori.}
Uno \textbf{stabilizzatore} (\textit{stabilizer})~\cite{gottesman1997} è un operatore $S$ che lascia invariati tutti gli stati validi del codice: $S\ket{\psi_L} = \ket{\psi_L}$. Finché lo stato resta nello spazio del codice, ogni misura di $S$ restituisce $+1$; un esito $-1$ segnala che un errore ha portato lo stato fuori dal sottospazio consentito. La collezione degli esiti di tutti i generatori del gruppo stabilizzatore costituisce la sindrome. Il punto cruciale è che gli stabilizzatori commutano con l'informazione logica: la loro misura proietta lo stato in un autospazio dell'errore, ma non distingue --- e quindi non collassa --- i coefficienti della sovrapposizione logica.

\paragraph{Distanza del codice.}
Un codice si denota $[[n, k, d]]$: $n$ qubit fisici, $k$ qubit logici, distanza $d$. La distanza è il peso minimo di un errore capace di trasformare uno stato logico valido in un altro senza essere rilevato; un codice di distanza $d$ corregge fino a
\begin{equation}
    t = \left\lfloor \frac{d-1}{2} \right\rfloor
    \label{eq:distanza_t}
\end{equation}
errori arbitrari. Lo Steane code ($d=3$) corregge un errore; un surface code di distanza 5 ne corregge due, di distanza 7 tre. La soglia oltre la quale aggiungere distanza \emph{conviene} --- il regime sotto-soglia in cui l'errore logico decresce all'aumentare di $d$ --- è l'oggetto quantitativo del Capitolo~\ref{chap:surface}, e la sua esistenza è garantita in linea di principio dal teorema della soglia per il calcolo fault-tolerant~\cite{aharonov1997fault}.

% -----------------------------------------------
\section{Il Panorama dei Codici e la Scelta Metodologica}
\label{sec:confronto_codici}

La progressione adottata è funzionale all'esperimento: il \emph{repetition code} rende leggibile l'estrazione di sindrome; lo Steane $[[7,1,3]]$ corregge un errore arbitrario; il surface code introduce distanza variabile, cicli rumorosi e un decoder scalabile. Le implementazioni dei primi due codici sono riportate nei Listati~\ref{lst:qec_repetition} e~\ref{lst:qec_steane} dell'Appendice~\ref{app:strumentiecodice}.

% -----------------------------------------------
\section{Laboratorio I: il Codice a Ripetizione}
\label{sec:repetition}

\subsection{Il codice bit-flip a 3 qubit}

Il codice a ripetizione contro errori X codifica $\ket{0_L} = \ket{000}$ e $\ket{1_L} = \ket{111}$: uno stato arbitrario $\alpha\ket{0} + \beta\ket{1}$ diventa $\alpha\ket{000} + \beta\ket{111}$ --- uno stato entangled, non tre copie. Gli stabilizzatori sono $Z_1 Z_2$ e $Z_2 Z_3$: misurano la \emph{parità} tra coppie di qubit senza misurare i singoli qubit. Un bit-flip su un qubit altera una o entrambe le parità, e la sindrome a 2 bit identifica univocamente quale dei tre qubit ha subito l'errore --- che viene corretto riapplicando X. La misura di parità avviene tramite ancilla: due CNOT copiano la parità sull'ancilla, che viene misurata al posto dei qubit dati.

Il limite è strutturale: il codice non rileva gli errori Z (phase-flip), che commutano con tutti gli stabilizzatori. Il duale --- codificare nella base di Hadamard, $\ket{\pm_L} = \ket{{\pm}{\pm}{\pm}}$, con stabilizzatori $X_1 X_2$ e $X_2 X_3$ --- protegge dai phase-flip ma non dai bit-flip. La dualità X/Z è la ragione per cui la correzione quantistica è più difficile di quella classica, e la concatenazione delle due protezioni è esattamente la costruzione del codice di Shor a 9 qubit~\cite{shorcorrection1995}.

\subsection{Risultati del laboratorio repetition}
\label{subsec:risultati_repetition}

Il primo controllo inietta deterministicamente un errore su ciascun qubit e verifica la lookup table
\[
00\!\to\! I,\qquad 10\!\to\! X_0,\qquad 11\!\to\! X_1,\qquad 01\!\to\! X_2.
\]
Le quattro sindromi osservate coincidono con la previsione; per coniugazione di Hadamard vale la stessa mappatura per gli errori $Z$ del codice duale. Il test completo e il circuito di estrazione sono documentati nel Listato~\ref{lst:qec_repetition}.

Il secondo test misura il beneficio della correzione. Si prepara lo stato logico, si applica a ciascun qubit dato un errore indipendente con probabilità $p$ (canale $X$ per il bit-flip, $Z$ per il phase-flip), si estrae la sindrome, si applica la correzione e si decodifica: la frazione di esiti logici errati su $20\,000$ ripetizioni (seed $=42$) stima l'errore logico $p_L$. La Figura~\ref{fig:qec_repetition} confronta i punti Monte Carlo con la previsione analitica del voto di maggioranza: il codice sbaglia quando almeno due dei tre qubit subiscono un errore, cioè
\begin{equation}
    p_L = 3p^2(1-p) + p^3 = 3p^2 - 2p^3.
    \label{eq:pL_repetition}
\end{equation}

\begin{figure}[H]
\centering
\includegraphics[width=0.75\textwidth]{qec_repetition_curve.pdf}
\caption[Errore logico del codice a ripetizione]{I punti Monte Carlo ($20\,000$ campioni per valore) seguono $p_L=3p^2-2p^3$; sotto $p=0.5$ la codifica riduce l'errore.}
\label{fig:qec_repetition}
\end{figure}

L'accordo è quantitativo: a $p=0.10$ si osserva $p_L=0.0288\pm0.0012$ contro $0.0280$ teorico. La soppressione è quadratica, ma il codice protegge un solo tipo di errore per volta; serve quindi un codice capace di correggere tutte le componenti di Pauli.

% -----------------------------------------------
\section{Laboratorio II: il Codice di Steane $[[7,1,3]]$}
\label{sec:steane}

\subsection{Struttura del codice}

Lo Steane code~\cite{steane1996} codifica un qubit logico in sette qubit fisici, ha distanza 3 e corregge un errore arbitrario (X, Z o Y) su un singolo qubit. È un codice \ac{CSS}, costruito a partire dal codice classico di Hamming $[7,4,3]$: la stessa matrice di parità classica genera sia gli stabilizzatori di tipo X sia quelli di tipo Z, il che permette di trattare separatamente le due famiglie di errori.

I generatori degli stabilizzatori, nella convenzione adottata, sono:
\begin{equation}
\begin{aligned}
\text{tipo X:} \quad & IIIXXXX, \quad IXXIIXX, \quad XIXIXIX \\
\text{tipo Z:} \quad & IIIZZZZ, \quad IZZIIZZ, \quad ZIZIZIZ
\end{aligned}
\label{eq:steane_stabilizzatori}
\end{equation}

Gli stabilizzatori di tipo Z rilevano gli errori X; quelli di tipo X rilevano gli errori Z; un errore Y, essendo proporzionale a $XZ$, produce entrambe le sindromi. La struttura di Hamming rende la decodifica trasparente: i tre bit di sindrome di ciascuna famiglia, letti come numero binario, indicano direttamente la \emph{posizione} del qubit colpito (sindrome $011 \rightarrow$ qubit 3, sindrome $110 \rightarrow$ qubit 6, e così via), e una lookup table sindrome~$\rightarrow$~correzione completa il decoder.

\subsection{Protocollo di laboratorio}

Il protocollo prepara $\ket{0_L}$, verifica la sindrome nulla, inietta $X$, $Z$ e $Y$ sui sette qubit e costruisce la lookup table; infine stima $p_L(p)$. L'estrazione di sindrome implementata non è fault-tolerant: il modello Monte Carlo applica gli errori ai qubit dati e assume la sindrome ideale. I risultati caratterizzano quindi il codice entro questo contratto, non un dispositivo protetto completo.

\subsection{Risultati del laboratorio Steane}
\label{subsec:risultati_steane}

Lo stato $\ket{0_L}$ è preparato con una rete di Hadamard e \ac{CNOT}; i sei stabilizzatori restituiscono sindrome nulla sui $2\,000$ campioni, verificando l'appartenenza al \textit{code space}. Il circuito completo è nel Listato~\ref{lst:qec_steane}.

\paragraph{La syndrome table è biiettiva.}
Il secondo controllo verifica i $21$ errori singoli $X$, $Z$ e $Y$. Le sette sindromi non nulle $001,\ldots,111$ identificano in modo biiettivo i sette qubit: gli stabilizzatori $Z$ diagnosticano $X$, quelli $X$ diagnosticano $Z$, mentre $Y$ attiva entrambe le famiglie. La tabella completa è quindi la rappresentazione binaria della posizione e non richiede un float dedicato.

\paragraph{Il qubit logico batte il fisico: soppressione quadratica.}
Sotto un canale depolarizzante di intensità $p$ per qubit, una simulazione Monte Carlo nel formalismo di Pauli ($200\,000$ campioni per punto) applica la correzione a peso minimo e verifica l'eventuale residuo logico. La Figura~\ref{fig:qec_steane} mostra $p_L\propto p^2$ con pendenza log--log $1.94$, come atteso per un codice di distanza 3 che fallisce a partire da errori di peso due.

\begin{figure}[H]
\centering
\includegraphics[width=0.72\textwidth]{qec_steane_curve.pdf}
\caption[Errore logico del codice di Steane {$[[7,1,3]]$}]{La pendenza $1.94$ conferma la soppressione quadratica; l'intersezione con $p_L=p$ dà una pseudo-soglia vicina a $0.08$.}
\label{fig:qec_steane}
\end{figure}

La pseudo-soglia è circa $p=0.08$: sotto tale valore il qubit logico è più affidabile del fisico. Rispetto al repetition code, Steane protegge simultaneamente da $X$, $Z$ e $Y$, ma resta a distanza fissa e l'estrazione qui simulata non include guasti delle ancille. Il surface code affronta questo limite mediante cicli ripetuti di sindrome rumorosa e distanza variabile.

\end{onehalfspace}
